\begin{Huge}
    Informatik (Bachelor of Education)-FAQ
\end{Huge}

\begin{block}{H\"aufig gestellte Fragen zum Studium}

\question{Lernt man im Studium, wie man programmiert?}{
    Ja, aber auf einer sehr eigenst\"andigen Basis. Man bekommt einen \"Uberblick \"uber die Sprache(n) und \"uber dahinterstehende Programmierkonzepte. Praktische Programmiererfahrung muss man sich aber gr\"o"stenteils selbst erarbeiten.
}

\question{Welche Programmiersprachen macht man da so?}{
    Das ist von Prof zu Prof und von Veranstaltung zu Veranstaltung unterschiedlich. In den ersten beiden Semestern lernt man meistens Racket und Java.
}

\question{Muss man programmieren k\"onnen, um das Studium anzufangen?}{
    Nein. Die Vorlesung beginnt absolut bei 0, um allen den Einstieg zu erm\"oglichen.
}

\question{Muss man gut in Mathe sein?}{
    Man muss kein Mathe-Genie sein, man sollte Mathe aber nicht hassen. Auch wenn man nur Mathe I hat, ben\"otigt man f\"ur die Informatik einiges an mathematischer Denkweise.
}

\question{Welche F\"acher kann man mit Informatik kombinieren?}{
    Alles au"ser NWT.
}

\question{Was, wenn ich merke, dass Lehramt doch nichts f\"ur mich ist?} {
    Du kannst Informatik gegen ein anderes Fach tauschen oder dich vom B.Ed. verabschieden und auf ein "`normales"' B.Sc.-Studium wechseln.
}

\question{Was mache ich denn mit einem Bachelor of Education?} {
    Die kurze Antwort: Weiter machen. Der Bachelor ist nur ein Teil deines Lehramtsstudiums, danach folgen noch der Master und das Referendariat.
}

\question{Wie sind meine Berufsschancen?} {
    Da in BaW\"u ein Mangel an Informatiklehrer*innen herrscht, sind deine Berufschancen sehr hoch.
}

\question{Was ist denn der Unterschied zwischen B.Sc. und B.Ed?} {
    Im B.Ed. h\"ort man ein paar weniger Informatikvorlesungen als im B.Sc. und hat daf\"ur Vorlesungen zu Bildungswissenschaften und Fachdidaktik.
}

\question{Was ist Fachdidaktik?} {
    In Fachdidaktik lernt man, wie man sein Wissen SuS~ (Sch\"ulerinnen und Sch\"ulern) beibringt und wie man Unterricht am besten planen kann.
}

\question{Gibt es Praktika?}{
    Ja, nachdem du BWS-I bestanden hast, musst du ein dreiw\"ochiges Orientierungspraktikum an der Schule machen, das in den Semesterferien stattfindet. Man hat zum Ablegen dieses Praktikums aber bis zum Ende des Bachelors Zeit. Im Master gibt es dann noch ein komplettes Semester Schulpraxis.
}

\question{Wie ist da so der NC?}{
    Gibt es keinen, der Studiengang ist zulassungsfrei.
}

\question{In welchen Klassenstufen wird Informatik unterrichtet?}{
    Hier ein kleines Schaubild zur Erl\"auterung, wie es eigentlich sein sollte:
    \begin{figure}[h!]
        \vspace{-1em}
        \includegraphics[width=0.6\linewidth]{images/info_bed_klassenstufen.png}
        \vspace{-1em}
    \end{figure}
    Momentan ist es aber leider so, dass sich Schulen aufgrund des Lehrermangels kaum an diesen Lehrplan halten k\"onnen.
}
\end{block}
